\thispagestyle{empty}

\noindent\rule[2pt]{\textwidth}{0.5pt}\\
{\textbf{Abstract ---}}
Probabilistic electricity demand forecasting is a critical component of modern
power system operations, enabling grid operators, energy traders, and capacity
planners to manage uncertainty arising from renewable integration, weather
variability, and evolving consumption patterns. This report presents a
systematic experimental study aimed at enhancing the AQ-NBEATS benchmark
model for probabilistic short-term electricity demand forecasting on the EMHIRES
dataset, covering 35 European countries over a 48-hour horizon. Starting from
the published baseline CRPS of 211.22, we design and evaluate more than twenty
experiments across four categories: exogenous feature and series-embedding
augmentation, curriculum quantile training, novel architectural contributions
(QCBC, TCR, DBE), and post-hoc Conformalized Quantile Regression (CQR).
Our best multi-seed result combines exogenous weather and calendar features
with adaptive quantile sampling, achieving a CRPS of 172.21 ,an 18.5\%
reduction over the published baseline, while CQR post-processing achieves
a nominal 95.96\% empirical coverage. These results establish new
state-of-the-art performance on the EMHIRES MHLV benchmark and provide
actionable guidelines for operational probabilistic load forecasting systems.\\
{\textbf{Keywords:}} Probabilistic forecasting, AQ-NBEATS, CRPS, EMHIRES,
conformal prediction, exogenous features, curriculum learning, electricity demand.\\
\noindent\rule[2pt]{\textwidth}{0.5pt}

\vspace{1.5cm}

\noindent\rule[2pt]{\textwidth}{0.5pt}\\
{\textbf{R\'esum\'e ---}}
La pr\'evision probabiliste de la demande d'\'electricit\'e est un composant
essentiel des op\'erations des syst\`emes \'electriques modernes. Ce rapport
pr\'esente une \'etude exp\'erimentale syst\'ematique visant \`a am\'eliorer
le mod\`ele de r\'ef\'erence AQ-NBEATS pour la pr\'evision probabiliste
de la demande \'electrique \`a court terme sur l'ensemble de donn\'ees EMHIRES,
couvrant 35 pays europ\'eens sur un horizon de 48 heures. En partant d'un CRPS
de r\'ef\'erence de 211.22, nous \'evaluons plus de vingt configurations
et atteignons un CRPS de 172.21, soit une r\'eduction de 18.5\%. La
calibration post-hoc par CQR atteint une couverture empirique de 95.96\%.\\
{\textbf{Mots cl\'es:}} Pr\'evision probabiliste, AQ-NBEATS, CRPS, EMHIRES,
pr\'ediction conforme, variables exog\`enes, apprentissage par curriculum.\\
\noindent\rule[2pt]{\textwidth}{0.5pt}
